%ctexart/ctexbook/ctexrep: 这些是ctex 宏包提供的文档类，适用于排版中文文章、书籍和报告
\documentclass[a4paper,UTF8]{ctexart}
%newtxtext 是一个 LaTeX 宏包，用于提供一个新的 Times 字体风格，适用于文本排版。它是 newtx 宏包系列的一部分，提供了一套新的 Times 风格字体，包括正文字体、粗体、斜体等，以及一些调整，用于更好地支持中英文混排。
%\usepackage{newtxtext}
\usepackage[left=2.5cm, right= 2.5cm, top=3.0cm, bottom=3.2cm]{geometry}
\usepackage[dvipsnames,svgnames]{xcolor}
\usepackage[strict]{changepage} % 提供一个 adjustwidth 环境
\usepackage{framed} % 实现方框效果
\definecolor{formalshade}{rgb}{0.95,0.95,1} % 文本框颜色
%%%%%%%%%%%%%%%%%%%%%%%%%%%%%%%%%%%%%%%%%%%%%%%%%%%%%%%%%%%%%%%%%%%%%%%%%%%%%%%%
\usepackage{titlesec}
\titleformat{\section}[block]{\normalfont\Large\bfseries}{\thesection}{1em}{}
\titlespacing*{\section}{0pt}{*3.5}{*1.5}
%%%%%%%%%%%%%%%%%%%%%%%%%%%%%%%%%%%%%%%%%%%%%%%%%%%%%%%%%%%%%%%%%%%%%%%%%%%%%%%%
\usepackage{listings}

\lstdefinestyle{customcpp}{
	language=C++,
	basicstyle=\small\ttfamily,
	keywordstyle=\color{blue},
	commentstyle=\color{green!50!black},
	stringstyle=\color{red},
	numbers=left,
	numberstyle=\tiny\color{gray},
	stepnumber=1,
	numbersep=10pt,
	backgroundcolor=\color{White!10},
	showspaces=false,
	showstringspaces=false,
	showtabs=false,
	frame=single,
	rulecolor=\color{White},
	tabsize=4,
	captionpos=b,
	breaklines=true,
	postbreak=\mbox{\textcolor{red}{$\hookrightarrow$}\space},
	escapeinside={(*}{*)},
	morekeywords={constexpr, nullptr, int, double, void, cout, endl} % 添加自定义关键字
}

\lstset{style=customcpp}

%%%%%%%%%%%%%%%%%%%%%%%%%%%%%%%%%%%%%%%%%%%%%%%%%%%%%%%%%%%%%%%%%%%%%%%%%%%%%%%%
\usepackage{fancyhdr}
% ------------------******-------------------
% 注意行末需要把空格注释掉，不然画出来的方框会有空白竖线
\newenvironment{formal}{%
\def\FrameCommand{%
\hspace{1pt}%
{\color{DarkBlue}\vrule width 2pt}%
{\color{formalshade}\vrule width 4pt}%
\colorbox{formalshade}%
}%
\MakeFramed{\advance\hsize-\width\FrameRestore}%
\noindent\hspace{-4.55pt}% disable indenting first paragraph
\begin{adjustwidth}{}{7pt}%
\vspace{2pt}\vspace{2pt}%
}
{%
\vspace{2pt}\end{adjustwidth}\endMakeFramed%
}
% ------------------******-------------------

 \usepackage{pdfpages}
\usepackage{graphicx}

%--------------------------------------------
\begin{document}
\includepdf[pages=-]{figures/background.pdf} % 
\setcounter{page}{1}%重置页码
	

\pagestyle{fancy}
\rhead{}
\lhead{}
\chead{Hello-algo}
\lfoot{}
\cfoot{-\ \thepage \  -}
\rfoot{}
%---------------------------------------------------------%
\section{初识算法}
\subsection{ 算法无处不在}
当我们听到“算法”这个词时，很自然地会想到数学。然而实际上，许多算法并不涉及复杂数学，而是更多
地依赖于基本逻辑，这些逻辑在我们的日常生活中处处可见。
\par 在正式探讨算法之前，有一个有趣的事实值得分享：你已经在不知不觉中学会了许多算法，并习惯将它们应
用到日常生活中了。下面，我将举几个具体例子来证实这一点。
\begin{formal}
C++中的仿函数（Functor）是一种特殊的对象，它可以像函数一样被调用，但实际上是一个类的实例。仿函数可以通过重载函数调用运算符 operator() 来实现，从而使得该对象可以像函数一样被调用。仿函数主要用于在算法中作为操作符使用，例如在STL（标准模板库）中的各种算法中。通过使用仿函数，可以将操作抽象出来，使得算法更加灵活，可以适用于不同的数据类型或者自定义的操作。
\end{formal}
\begin{lstlisting}[language=C++, label=lst:cpp]
#include <iostream>

// 定义一个加法仿函数
class AddFunctor {
public:
    int operator()(int a, int b) {
        return a + b;
    }
};

int main() {
    AddFunctor adder;  // 创建仿函数对象

    int result = adder(3, 5);  // 使用仿函数调用
    std::cout << "Result:" << result << std::endl;

    return 0;
}
\end{lstlisting}


\end{document}
